\documentclass[11pt]{article}
\newcommand{\SheetCode}{Theory 04}
\usepackage[margin=25mm]{geometry}
\usepackage[T1]{fontenc}
\usepackage{lmodern}
\DeclareUnicodeCharacter{2605}{$\star$}
\usepackage{microtype}
\usepackage{amsmath,amssymb,mathtools,bm}
\usepackage{enumitem}
\usepackage{xcolor}
\usepackage{fancyhdr}
\usepackage{hyperref}
\usepackage[most]{tcolorbox}
\definecolor{agiBlue}{HTML}{173B57}
\definecolor{agiGold}{HTML}{B7791F}
\definecolor{agiLight}{HTML}{EFF6FA}
\definecolor{agiGreen}{HTML}{184D3B}
\hypersetup{colorlinks=true,linkcolor=agiBlue,urlcolor=agiBlue}
\setlength{\parindent}{0pt}
\setlength{\parskip}{0.55em}
\setlength{\headheight}{14pt}
\setlist{nosep,leftmargin=*}
\pagestyle{fancy}
\fancyhf{}
\fancyhead[L]{\small Part of the \href{https://github.com/mmikhasenko/GIModel.jl}{Agentic Gottried--Isgur} project}
\fancyhead[R]{\SheetCode}
\fancyfoot[C]{\thepage}
\newcommand{\dd}{\,\mathrm d}
\newcommand{\ket}[1]{\lvert #1\rangle}
\newcommand{\bra}[1]{\langle #1\rvert}
\newcommand{\braket}[2]{\langle #1\mid #2\rangle}
\newcommand{\expect}[1]{\left\langle #1\right\rangle}
\newcommand{\op}[1]{\widehat{#1}}
\newcommand{\GeV}{\,\mathrm{GeV}}
\newcommand{\R}{\mathbb{R}}
\newtcolorbox{goals}{colback=agiLight,colframe=agiBlue,title=Learning outcomes}
\newtcolorbox{reading}{colback=white,colframe=agiGold,title=Book reference,
  after skip=5mm}
\newtcolorbox{paperbridge}{colback=agiGreen!7,colframe=agiGreen,
  colbacktitle=agiGreen,coltitle=white,title=Connection to the GI paper}
\newtcolorbox{problem}[1]{colback=white,colframe=agiBlue,title={#1},
  before skip=3.5mm,after skip=1mm}
\newtcolorbox{solution}{colback=black!2,colframe=black!45,title=Solution}
\newtcolorbox{quiz}{colback=white,colframe=agiGold,title=Post-solution concept check}
\newtcolorbox{checkpoint}{colback=white,colframe=agiGold,title=Interpretation checkpoint}
\newcommand{\sheettitle}[3]{%
  \begin{center}
  {\Large\bfseries #1}\par
  \vspace{0.2em}{\color{agiBlue}#2}\par
  \vspace{0.4em}\small #3
  \end{center}\vspace{0.5em}}
\newcommand{\difficulty}[1]{\textbf{Difficulty:} #1}
\newcommand{\timebox}[1]{\textbf{Suggested time:} #1}
\newcommand{\answer}{\par\smallskip\color{black!50}\textbf{Answer:} }

\begin{document}
\sheettitle{Sheet 4: The spin-independent GI interaction}{Color, confinement, and a running Coulomb term}{Prerequisite: Sheets 1--3. \difficulty{★/★★★} \quad \timebox{2--3 hours}}

\begin{goals}
Reduce the color interaction in a singlet meson; understand the Cornell-like
central potential; transform the Gaussian running coupling; analyze its
short- and long-distance limits; predict qualitative flavor trends.
\end{goals}
\begin{reading}
Selected reading in Landau--Lifshitz, \emph{Quantum Mechanics: Non-Relativistic
Theory}: Ch. III, ``Schr\"odinger's Equation,'' and Ch. V, ``Motion in a
Centrally Symmetric Field,'' especially the Coulomb problem. This is a useful
comparison, not a prerequisite or the structure followed by the sheet.
\end{reading}


For a color-singlet $q\bar q$ pair,
$\expect{\bm F_1\cdot\bm F_2}=-4/3$. The reduced central form is
\[
V(r)=br+c-\frac{4\alpha_s(r)}{3r},\qquad
\alpha_s(r)=\sum_k a_k\operatorname{erf}(\gamma_k r).
\]

\begin{problem}{Problem 1 ★ — Recover the singlet potential}
The color interaction is
$V_{\rm color}=-[\tfrac34c+\tfrac34br-\alpha_s(r)/r]\bm F_1\cdot\bm F_2$.
For a singlet, $\expect{\bm F_1\cdot\bm F_2}=-4/3$.
Compute $V(r)=\expect{V_{\rm color}}$, then identify its constant,
confining, and Coulomb pieces. Take $b>0$ and $\alpha_s(r)>0$;
state which piece is attractive and which grows at large separation.
\end{problem}
\begin{solution}
Multiplication by $-(-4/3)=4/3$ gives
$c+br-4\alpha_s(r)/(3r)$. The attractive Coulomb sign and rising confinement
term follow simultaneously from the singlet color projection.
\end{solution}
\begin{quiz}
Which part of the potential prevents an isolated quark pair from separating at
finite energy? \answer the linearly rising $br$ term.
\end{quiz}

\begin{problem}{Problem 2 ★★ — Derive the error function}
Use
$\int \dd^3q\,(2\pi)^{-3}4\pi e^{-q^2/(4\gamma^2)}e^{i\bm q\cdot\bm r}/q^2$
to show that a Gaussian-regulated Coulomb kernel is
$\operatorname{erf}(\gamma r)/r$. You may differentiate the radial integral
with respect to $r$.

\textbf{Hint:} Take $\gamma>0$ and the polar axis along $\bm r$.
Use $\int\dd\Omega_q\,e^{iqr\cos\theta}=4\pi\sin(qr)/(qr)$,
$\operatorname{erf}(x)=2\pi^{-1/2}\int_0^x e^{-t^2}\dd t$, and
$\int_0^\infty e^{-aq^2}\cos(qr)\dd q=\sqrt\pi e^{-r^2/(4a)}/(2\sqrt a)$.
Differentiate $I(r)=\int_0^\infty e^{-q^2/(4\gamma^2)}\sin(qr)\dd q/q$,
then integrate back using $I(0)=0$.
\end{problem}
\begin{solution}
Angular integration reduces the transform to
$\frac{2}{\pi r}\int_0^\infty e^{-q^2/(4\gamma^2)}\sin(qr)/q\,\dd q$.
Differentiating its integral with respect to $r$ produces the elementary cosine
Gaussian integral; integrating back from zero gives
$(\pi/2)\operatorname{erf}(\gamma r)$, leaving
$\operatorname{erf}(\gamma r)/r$.
\end{solution}
\begin{quiz}
Why is the coordinate-space answer finite at $r=0$? \answer
$\operatorname{erf}(x)\sim2x/\sqrt\pi$, cancelling the explicit $1/r$.
\end{quiz}

\begin{problem}{Problem 3 ★ — Compute both distance limits}
Find the $r\to0$ and $r\to\infty$ limits of $\alpha_s(r)$ and of the Coulomb
piece $G(r)=-4\alpha_s(r)/(3r)$.

\textbf{Hint:} Assume finite $a_k$ and $\gamma_k>0$.
Use $\operatorname{erf}(x)=2x/\sqrt\pi+O(x^3)$ near zero and
$\operatorname{erf}(x)\to1$ as $x\to+\infty$. Give the leading behavior
as well as the limiting value of $G$.
\end{problem}
\begin{solution}
At small $r$,
$\alpha_s(r)=(2r/\sqrt\pi)\sum_k a_k\gamma_k+O(r^3)$, so
$G(0)=-8\sum_k a_k\gamma_k/(3\sqrt\pi)$ is finite. At large $r$,
$\alpha_s\to\sum_k a_k\equiv\alpha_s^{\rm crit}$ and
$G\sim-4\alpha_s^{\rm crit}/(3r)$. With GI coefficients $0.25,0.15,0.20$,
$\alpha_s^{\rm crit}=0.60$.
\end{solution}
\begin{quiz}
Does “freezing” mean the potential becomes constant at large $r$? \answer no;
the coupling freezes while its Coulomb potential still falls as $1/r$.
\end{quiz}

\begin{problem}{Problem 4 ★★ — Derive a virial scale}
Model the energy of a state of size $R>0$ by
$E(R)=1/(2\mu R^2)+bR-\kappa/R$, with $\mu,b,\kappa>0$.
Here $R$ is a trial size, not a radial wavefunction. Find the stationary
equation and show it has exactly one positive solution, which is a minimum.
Holding $b,\kappa$ fixed, determine whether that size increases or decreases
with $\mu$. No explicit cubic-root solution is required.

\textbf{Hint:} Multiply $E'(R)=0$ by $R^3$, then differentiate implicitly.
\end{problem}
\begin{solution}
$E'(R)=-1/(\mu R^3)+b+\kappa/R^2=0$, or
$bR^3+\kappa R=1/\mu$. The function $bR^3+\kappa R$ increases strictly from zero to infinity,
so there is one positive solution. The sign of $E'$ changes there from
negative to positive, giving a minimum. Implicit differentiation gives
$dR/d\mu=-1/[\mu^2(3bR^2+\kappa)]<0$: the estimated state shrinks as
its reduced mass increases.
\end{solution}
\begin{quiz}
Why does compactness make heavy quarkonium a clean test of the model? \answer
the system is more nonrelativistic and less sensitive to long-distance light
quark dynamics, while still testing the short-distance potential.
\end{quiz}

\begin{problem}{Problem 5 ★★★ — Differentiate the Coulomb profile used by GI}
Derive $G'(r)$ for $G=-4\alpha_s(r)/(3r)$ and
$\alpha_s'=\sum_k(2a_k\gamma_k/\sqrt\pi)e^{-\gamma_k^2r^2}$.
State why differentiating the analytic profile is preferable to differentiating
tabulated values.
\end{problem}
\begin{solution}
\[
G'(r)=\frac{4\alpha_s(r)}{3r^2}-\frac{4\alpha_s'(r)}{3r}.
\]
Analytic differentiation preserves the same running-coupling formula and avoids
small point-to-point numerical errors that derivatives can amplify, especially
when followed by division by $r$. A finite-difference derivative can then serve
as an independent check of the analytic result rather than becoming part of the
definition of the interaction.
\end{solution}
\begin{quiz}
If an analytic derivative and a finite-difference check disagree only very near
$r=0$, what should be examined first? \answer use the series limit in which the
apparently singular terms cancel, and then vary the finite-difference step.
Changing a physical model parameter would not address either numerical issue.
\end{quiz}

\begin{paperbridge}
GI Eq. (3) gives the linear-plus-Coulomb confinement interaction and Eq. (9)
gives the meson color factor $-4/3$. Eqs. (12)--(13) relate the Gaussian sum for
$\alpha_s(Q^2)$ to the coordinate-space error-function form; the numerical
coefficients are printed in the caption of Fig. 2.
\end{paperbridge}
\end{document}
